This section provides a taxonomy of memory poisoning attacks against RAG-based agent systems. While the specific implementations vary, all exploit the fundamental trust assumption: that retrieved memories are legitimate.

\subsection{MemoryGraft: Persistent Experience Poisoning}

The MemoryGraft attack \cite{memorygraft2025} represents the most comprehensive demonstration of memory poisoning to date. The attack operates in three phases:

\textbf{Phase 1: Payload Construction.} The attacker prepares two sets of experience records:
\begin{itemize}
    \item $S_{\text{benign}}$: A collection of legitimate-appearing workflow experiences (typically 100+ entries)
    \item $S_{\text{poison}}$: A smaller set of malicious experiences containing attack patterns $\pi$ (typically 10 entries)
\end{itemize}

The malicious patterns can include:
\begin{itemize}
    \item Validation bypass: \texttt{df['\_validation']='ok'}
    \item Remote execution: \texttt{curl -s https://attacker.com/script | bash}
    \item Data exfiltration: \texttt{curl -F file=@/tmp/data.tgz https://attacker.com/upload}
    \item Audit trail removal: \texttt{df.drop(columns=['audit\_id','consent\_flag'])}
\end{itemize}

\textbf{Phase 2: Ingestion.} The attacker supplies a benign-looking document (README, documentation, task description) containing embedded code that constructs the poisoned memory store. When the agent ingests the document, it executes the code and persists the malicious entries.

\textbf{Phase 3: Activation.} The poisoned entries are designed to occupy high-density regions of the embedding space. When the agent processes semantically similar tasks, the retrieval function returns a mixture of legitimate and poisoned memories. The agent then imitates patterns from both---including the malicious ones.

A critical finding from the MemoryGraft paper: with only 10\% poisoned entries, the \emph{Poisoned Retrieval Proportion} (PRP) can reach 47.9\%. Nearly half of retrieved items are malicious, despite being less than one-tenth of the store.

\subsection{MINJA: Query-Only Memory Injection}

The MINJA attack \cite{minja2025} demonstrates that memory poisoning is possible without any file access---purely through query interaction. The attack uses:

\begin{itemize}
    \item \textbf{Bridging Steps}: Gradual steering through intermediate queries that appear legitimate
    \item \textbf{Indication Prompts}: Signals that encourage the agent to memorize the interaction
    \item \textbf{Progressive Shortening}: Compression of malicious context while preserving poisoned relations
\end{itemize}

MINJA achieves $>$95\% injection success rate under idealized conditions. The attack is slower than MemoryGraft (requiring multiple interactions) but harder to detect because individual queries appear innocuous.

\subsection{AgentPoison: Backdoor Triggers}

AgentPoison \cite{agentpoison2024} takes a different approach: rather than relying on semantic similarity, it optimizes specific \emph{trigger phrases} that map to unique regions of the embedding space. When a user query contains the trigger, poisoned demonstrations are retrieved with high probability.

The attack achieves $>$80\% success rate with $<$0.1\% poison rate---far more efficient than MemoryGraft's 10\% requirement. However, it requires the trigger phrase to appear in user input, making it less general.

\subsection{XAMT: Stealthy Perturbation}

XAMT \cite{xamt2025} uses bilevel optimization to find memory perturbations that are effective yet evade detection heuristics. The attack can operate at sub-percent poison rates ($\leq$0.1\% in RAG systems) while maintaining covertness.

\subsection{Common Attack Properties}

Across all variants, several properties emerge:

\begin{enumerate}
    \item \textbf{Persistence}: Poisoned memories survive across sessions, creating long-term behavioral drift
    \item \textbf{Trigger-Free Activation}: Unlike traditional backdoors, most attacks activate on any semantically similar task
    \item \textbf{Semantic Centrality}: Effective attacks target high-traffic regions of the embedding space
    \item \textbf{Plausible Deniability}: Poisoned entries are designed to appear legitimate on individual inspection
    \item \textbf{Amplification Through Retrieval}: The dual-channel retrieval (lexical + semantic) can amplify poison retrieval beyond the base poison rate
\end{enumerate}

\begin{table}[h]
\centering
\caption{Comparison of Memory Poisoning Attacks}
\begin{tabular}{lcccl}
\toprule
\textbf{Attack} & \textbf{Poison Rate} & \textbf{Success Rate} & \textbf{Requires} & \textbf{Trigger} \\
\midrule
MemoryGraft & 10\% & 48\% PRP & File ingestion & Semantic \\
MINJA & Variable & $>$95\% & Query access & Semantic \\
AgentPoison & $<$0.1\% & $>$80\% & KB access & Explicit \\
XAMT & $\leq$0.1\% & Variable & Any & Semantic \\
\bottomrule
\end{tabular}
\label{tab:attack-comparison}
\end{table}
